\section{Statements and notation from \texorpdfstring{\cite{OAI26}}{[OAI26]}}\label{app:oai}

This appendix records the symbol translation between~\cite{OAI26} (version of 23 September 2026) and the present paper,
and restates the results of~\cite{OAI26} cited as black boxes in Section~\ref{sec:prelim-oai}
(Facts~\ref{fact:oai-streams}--\ref{fact:oai-networks}) together with the printed form of~\cite[Prop.~14]{OAI26} for
comparison with Proposition~\ref{prop:prop14-param}.

\paragraph{Notation map.} Table~\ref{tab:oai-renaming} lists every symbol of~\cite{OAI26} that is renamed here to avoid
collision with the bridge parameters. Headers that concatenate positive integer parameters use the self-delimiting
code~\cite[(1)]{OAI26}:
\begin{equation}\label{eq:oai-gamma}
  b_v=\lfloor\log_2v\rfloor+1,\qquad \Gamma(v)=0^{\,b_v-1}\,\mathrm{bin}_{b_v}(v),\qquad |\Gamma(v)|=2b_v-1
  \qquad(v\ge1).
\end{equation}

\begin{table}[htbp]
\centering\small
\caption{Symbol correspondence between~\cite{OAI26} and this paper.}\label{tab:oai-renaming}
\begin{tabular}{@{}>{\raggedright}p{4.0cm}>{\raggedright}p{3.4cm}>{\raggedright\arraybackslash}p{8.1cm}@{}}
\toprule
\cite{OAI26} & This paper & Meaning\\
\midrule
$P$ (\S2, Lemmas~3, 12, Prop.~14) & $A$ & collapsed prefix address range ($P$ is our bridge working precision)\\
$q$ (\S2, Lemmas~3--4; \S4, Lemmas~10--12) & $q_{\mathrm{rad}}$ & fixed integer or odd prime radix ($q$ is our layer grid precision)\\
$p$ (\S\S5--9, Prop.~14) & $q$ & dyadic grid precision of layers and middle transforms ($q=P+2$ in \S\ref{sec:bridge})\\
$R$ (\S2, Lemma~3; \S5.2, Lemma~13) & $R_{\mathrm{row}},\,R_{\mathrm{rec}}$ & row-field range; record bit width ($R=t_1\cdots t_d$ is our DFT length)\\
$E,\,B,\,\nu,\,k$ (\S5.4) & $E_{\mathrm{net}},\,B_{\mathrm{net}},\,\nu_{\mathrm{net}},\,j_{\max}$ & complex-network guard constants, guard degree, and recursion depth ($k=\lg R$)\\
$T,\,L=\log_2T$ (\S\S6--9) & $R,\,k=\lg R$ & total power-of-two transform box size and its binary logarithm\\
$\epsilon,\,\log$ & $\eps,\,\ln$ & dimension exponent; natural logarithm ($\lg=\log_2$)\\
\bottomrule
\end{tabular}
\end{table}

\subsection{Machine model and elementary stream lemmas}\label{sec:prelim-oai-model}

In~\cite[\S2]{OAI26}, a multitape Turing machine has a fixed finite alphabet and a fixed finite number of
one-dimensional tapes; an array on $[a_1]\times\dots\times[a_t]$ is stored in lexicographic order (first coordinate
most significant) with implicit addresses, logical bit volume $V$, and move convention $(fA)(f(x))=A(x)$. The following
two lemmas of~\cite[\S2.1]{OAI26} constitute Fact~\ref{fact:oai-streams}.

\begin{proposition}[{\cite[Lemma~3, \S2.1]{OAI26}}]\label{oai:lem3}
Fix a radix $q_{\mathrm{rad}}\ge2$. On an array of volume $V$, each of the following takes $O(V)$ steps on a fixed
number of one-dimensional tapes (including shape preparation, counter initialization/resets, and workspace cleanup):
(1)~any fixed permutation of the pair of digit values at two specified radix-$q_{\mathrm{rad}}$ positions;
(2)~moving one specified radix-$q_{\mathrm{rad}}$ digit to another position, preserving the order of intervening
positions;
(3)~splitting on a specified radix-$q_{\mathrm{rad}}$ digit, applying a fixed pointwise map to the aligned streams, and
merging; and
(4)~within each prefix of an array on $[A]\times[R_{\mathrm{row}}]\times[B]$, splitting whole rows cyclically into a
fixed number $W$ of streams and, when $W\mid R_{\mathrm{row}}$, merging them in any fixed order of the row classes.
\end{proposition}

\begin{proposition}[{\cite[Lemma~4, \S2.1]{OAI26}}]\label{oai:lem4}
Fix a prime $q_{\mathrm{rad}}$, a fixed number of control fields, and finitely many rational coefficients with
denominators prime to $q_{\mathrm{rad}}$ (and numerators prime to $q_{\mathrm{rad}}$ when used as units). Put
$Q=q_{\mathrm{rad}}^b$ ($b\ge1$), and let a target field $y\in[Q]$ follow its control fields, each control used in an
offset having $O(b)$ digits, with arbitrary spectator fields and suffix records. An update
$y\mapsto y+a(\text{controls})\pmod Q$ costs $O(V)$ whenever the offset is computable in $b^{O(1)}$ time from the
controls (in particular for fixed rational linear combinations modulo $Q$), and so does multiplication of $y$ by any of
the specified rational units or its inverse. Fixed-width signed fixed-point addition, subtraction, sign change, and
shifts by fixed powers of two (with sufficient guard bits), as well as intra-record bit reversal, also have $O(V)$ bit
cost.
\end{proposition}

\subsection{Finite-network certificates and printed butterfly layer}

At $h=100$, $v=\binom{h}{3}=161\,700$, $N=v^3$, and $m=h^3=10^6$, \cite[\S3.1]{OAI26} constructs a three-stage
bit network over $\mathbb F_2$ ($W_b$ wires, neighbor rule $|S\cap T|=1$) and complex network over $\Z[i,1/2]$ ($W_c$
wires, distinct triples with $|S\cap T|$ even) on two data banks $X_a,Y_a$ ($a\in\binom{[h]}{3}^3$) and per-invocation
scratch wires. By~\cite[Lemmas~5--6, \S\S3.1--3.3]{OAI26}, the forward-inverse-forward schedule sends
$(X_a,Y_a)\mapsto(-Y_a,X_a)$, restores all scratch wires from arbitrary initial values, and attaches nested
nondegenerate subspaces of $\mathcal F=F^{\otimes3}$ ($\dim\mathcal F=m$) whose only decreasing edges lose dimension $h$
on central wires ($L_b/N=100/539$, $L_c/N=101/539$). Here $F=\Q^h$ under $x^{\mathsf T}(I_h-J_h/9)y$ for the bit
network (where $I_h$ and $J_h$ are the $h\times h$ identity and all-ones matrices), and $F=\mathbb F_2^h$ under the dot
product for the complex network. The resulting
certificates form Fact~\ref{fact:oai-networks}:

\begin{proposition}[{\cite[Props.~7--8, \S\S3.4--3.5]{OAI26}}]\label{oai:networks}\label{oai:prop8}
\begin{enumerate}[label=(\alph*),leftmargin=*]
\item \emph{Bit network (\cite[Prop.~7]{OAI26}).} Assigning orthogonal projections $P_U$ at gates/sinks, $-P_{U_a}$ at
  sources $X_a$ and $0$ at other sources, the bit network uses only pointwise XOR gates on $W_b$ roles, induces the role
  permutation $\rho_b$ exchanging $X_a$ and $Y_a$ while fixing scratch roles, satisfies
  $M_{\mathrm{sink}(\rho_b(w))}-M_{\mathrm{source}(w)}=I_m$ for every input role $w$, and has total rational edge-rank
  sum $s_b=W_bm-N+2L_b<W_bm$ with the exact integers of Definition~\ref{def:fixed-network-premises}.
\item \emph{Complex network (\cite[Prop.~8]{OAI26}).} With $H_{\mathrm u}=\frac1{\sqrt2}\bigl(\begin{smallmatrix}1&1\\1&-1\end{smallmatrix}\bigr)$,
  $C=H_{\mathrm u}\diag(1,i)H_{\mathrm u}=aI+bX$ ($a=\frac{1+i}2$, $b=\frac{1-i}2$), $q_U(x)=\operatorname{wt}(P_Ux)\bmod4$
  over $\mathbb F_2$ and $C_U=H_{\mathrm u}^{\otimes m}\diag_x(i^{q_U(x)})H_{\mathrm u}^{\otimes m}$, each edge frame
  difference $C_VC_U^{-1}$ is a product of one translation kernel $aI+bX_v$ or its inverse per vector in an orthonormal
  basis of the edge residual, with total basis length $s_c=W_cm-2N+2L_c<W_cm$ as in
  Definition~\ref{def:fixed-network-premises}. After source/sink diagonal sign corrections ($Z_{a,f}$ before $X_a$ and
  $i^{27f}Z_{a,f}$ after $Y_a$ on $f$ columns) and undoing the signed bank exchange, telescoping yields
  $(C^{\otimes m})^{\otimes f}$ on every role, including scratch roles with arbitrary initial values.
\end{enumerate}
\end{proposition}

By~\cite[(18), \S5.1]{OAI26}, $C^{\otimes D}$ converts to the normalized butterfly layer $H_0^{\otimes D}$ via
\begin{equation}\label{eq:oai-h0}
  H_0^{\otimes D}=b^D\,\mathcal S_D\,C^{\otimes D}\,\mathcal S_D,\qquad
  (\mathcal S_Dz)(x)=i^{x_1+\dots+x_D}z(x),\qquad b^D=(-i/2)^{\lfloor D/2\rfloor}b^{D\bmod2}.
\end{equation}
Finally, for comparison with Proposition~\ref{prop:prop14-param}, we record~\cite[Prop.~14]{OAI26} as printed:

\begin{proposition}[{\cite[Prop.~14, \S5.4]{OAI26}}, printed constants]\label{oai:prop14}
 Proposition~\ref{prop:prop14-param} holds as stated in~\cite[Prop.~14]{OAI26} at the fixed structural constants
$\beta=\frac12$, $\tau=\sigma=1-2^{-50}$ and $C_1=20$ (with $q$ written $p$ and $A$ written $P$ there), for any fixed
positive rationals $c,\eps,\lambda,\lambda'$ satisfying~\eqref{eq:walsh-layer-a}--\eqref{eq:walsh-layer-b} and
$\eps C_1<1$.
\end{proposition}
